%%%%%%%%%%%%%%%%%%%%%%preface.tex%%%%%%%%%%%%%%%%%%%%%%%%%%%%%%%%%%%%%%%%%
% sample preface
%
% Use this file as a template for your own input.
%
%%%%%%%%%%%%%%%%%%%%%%%% Springer %%%%%%%%%%%%%%%%%%%%%%%%%%

\preface

%% Please write your preface here
% Use the template \emph{preface.tex} together with the Springer document class SVMono (monograph-type books) or SVMult (edited books) to style your preface in the Springer layout.

% A preface\index{preface} is a book's preliminary statement, usually written by the \textit{author or editor} of a work, which states its origin, scope, purpose, plan, and intended audience, and which sometimes includes afterthoughts and acknowledgments of assistance. 

% When written by a person other than the author, it is called a foreword. The preface or foreword is distinct from the introduction, which deals with the subject of the work.

% Customarily \textit{acknowledgments} are included as last part of the preface.

\paragraph{Origin}
In this book we present series of tutorials and advanced topics originally written as documentation parts for the \Digraph collection of Python ressources over the last ten years \citep{BIS-2021}. These \Digraph resources -- like the \texttt{outrankingDigraphs} module -- were essentially used for the preparation and illustration of the \emph{Algoritmic Decision Theory} Course taught at the University of Luxembourg from 2010-2020. Some parts also served for preparing and illustrating the Lectures of a \emph{Computational Statistics} Course like the \texttt{randomNumbers} module. Curious readers will also discover some resources like the \texttt{arithmetics} module, used for preparing and illustrating a first Semester Course on Discrete Mathematics.


\paragraph{Scope}

The \Digraph documentation, available \href{https://digraph3.readthedocs.io/en/latest/}{on the Read The Docs site}, describes the Python3 resources for implementing decision aid algorithms in the context of a bipolar-valued outranking approach. These computing resources are useful in the field of Algorithmic Decision Theory and more specifically in outranking based Multiple Criteria Decision Aid (MCDA).

We address essentially three types of usage. First, we provide present methods an illustrate tools for solving either a multiple criteria best choice selection or a dual worst choice rejection problem, mainly interesting for management or political studies. We also address the problem of how to list a set of items with multiple incommensurable performance criteria either from the best to the worst (ranking problem) or from the worst to the best (ordering problem), mainly interesting for production schedulers or recommender systems designers. The third type of usage should mainly be of interest for performance auditors, as we present methods and tools for relative or absolute quantiles rating of mulitiple criteria performance records. 

The \Digraph ressources do not provide a professional Python software library. The modules were not built following any professional software development methodology. The design of classes and methods was kept as simple and elementary as was opportune for the author. Sophisticated and criptic overloading of variables, methods and classes is more or less avoided all over. A simple copy, paste and ad hoc customization development strategy is generally preferred. As a consequence, the modules keep a large part of independence and are easier to maintain and adapt.  Furthermore, the development of the \Digraph modules being spread over a decade, our programming style did evolve with our experience and the changes and enhancement of the standard Python libraries. Backward comptbility necessarily introduced so with time some notation and programming technique changes.

Finally, the book is not presenting any mathematical developments and proofs. Readers interested in the mathematical aspects of our methods and tools are invited to consult the references. Full texts of most references may be downloaded from the \href{https://orbilu.uni.lu/}{https://orbilu.uni.lu/} repository of the University of Luxembourg. 

\paragraph{Plan}

\noindent The book contains, first, a set of chapters introducing the \Digraph resources and the main formal objects discussed in this book, like digraphs and outranking digraphs. A second, methodological part presents evaluation models and decision methods and tools. Some chapters are problem oriented and show how to compute the winner of an election, how to build a best choice recommendation, or how to linearly rank or rate with multiple incommensurable performance criteria. A third part present case studies for best choice, for ranking and for rating with multiple incomensurabel performance criteria. A fourth Section presents more advanced topics showing some pearls of bipolar-valued epistemic logic. We generalize Kendall's ordinal correlation index to bipolar-valued digraphs in order measure fitness of ranking results. A major Chapter shows how to compute digraph kernels and solve Berge kernel equaltion systems. We furthermore introduce confidence levels for outranking situations with incertain significance weights. We also develop a robustness analysis of the outranking digraph when facing ordinal criterion significance weights. A last part concerns simple undirected graphs and present operational aspects of computing maximal independent sets (MISs) and kernels in graphs and digraphs. The tutorial about split, interval and permutation graphs is inspired by Martin Golumbic ‘s book on Algorithmic Graph Theory and Perfect Graphs. We also provide a tutorial on tree graphs and spanning forests.

\paragraph{Intended audience}

Master students and doctoral candidates in Computer Science, Engineering Sciencess and Computational Management Sciences taking a course on Algorithmic Decision Theory, Multiple criteria Decision Aid or DEcision Analysis.

\paragraph{Aknowledgments}

This book contains many ideas, methods and tools that are not only the author’s. They have been borrowed from friends and colleagues: 

\vspace{0.5cm}
\begin{minipage}{7cm}
\emph{Denis Bouyssou, Luis Dias,}\\ 
\emph{Claude Lamboray, Patrick Meyer,}\\
\emph{Vincent Mousseau, Alex Olteanu,}\\
\emph{Marc Pirlot, late Bernard Roy,}\\
\emph{Alexis Tsouki\`as, Thomas Veneziano,}\\
and especially,\\
\emph{late Marc Roubens}.
\end{minipage}\quad
\begin{minipage}{3cm}
\includegraphics[width=3cm]{Figures/Marc-Roubens.jpg} \\
{\tiny \emph{Marc Roubens}}
\end{minipage}

\vspace{0.3cm}
Their help is gratefully acknowledged.

\vspace{\baselineskip}
\begin{flushright}\noindent
Luxembourg, 2021\hfill {\it Raymond Bisdorff}\\
%month year\hfill {\it Firstname  Surname}\\
\end{flushright}


